Dado el siguiente problema de programación lineal:
\begin{align*}
\text{max.} \quad & z = 4x_1 + x_2 \\
\text{s.a.} \quad & 3x_1 + 2x_2 \leq 6 \\
                 & 6x_1 + 3x_2 \leq 10 \\
                 & x_1, x_2 \geq 0
\end{align*}
Se ha obtenido una tabla correspondiente a la matriz completa cuya fila 0 es la siguiente:

\begin{table}[h]
\begin{center}
\begin{tabular}{ccccc}
$z$                         & $x_1$ & $x_2$ & $h_1$ & $h_2$                   \\
\multicolumn{1}{|l|}{-20/3} & 0     & -2    & 0     & \multicolumn{1}{l|}{-1} \\ \hline
\end{tabular}
\end{center}
\end{table}
    

Demuestra, \textbf{\textit{sin resolver el problema y utilizando la teoría de la dualidad}}, que tiene que haber algún cálculo erróneo en la fila 0.



El problema dual asociado es:
\begin{align*}
\min \quad & w = 6y_1 + 10y_2 \\
\text{s.a.} \quad & 3y_1 + 6y_2 \geq 4 \\
                  & 2y_1 + 3y_2 \geq 1 \\
                  & y_1, y_2 \geq 0
\end{align*}

Si se ha llegado a la solución óptima del primal, el valor de la función objetivo del dual tendría que ser $w^*=20/3$ según el teorema fundamental de la dualidad y las variables en la base del dual serían $y_2=1$ y $h'_2=2$.

Esta solución del dual no es factible (no se cumple la primera restricción del dual), luego tiene que haber un error en el cálculo de la fila 0 de la matriz completa de la solución óptima del problema inicial.

Además, $z^*_P=\frac{20}{3}\neq 10=s^*_D$.

